\clearpage
\phantomsection
\addcontentsline{toc}{chapter}{Tóm tắt}
\chapter*{\fontsize{13}{15}\selectfont{Tóm tắt}}
\fontsize{12}{12}\selectfont{
\noindent\textbf{Tóm tắt:}
Dự đoán cảm xúc trong hội thoại là bài toán ước lượng một người sẽ cảm thấy thế nào ở lượt nói tiếp theo, trước khi lượt nói đó diễn ra. Bộ dữ liệu Hi-EF đặt bài toán này trên video phim truyền hình: từ ba đoạn hội thoại liên tiếp, mô hình phải dự đoán cảm xúc của người sẽ nói ở đoạn thứ tư, trong khi danh tính người đó không được cho trước. Báo cáo trình bày quá trình nghiên cứu bài toán, từ phân tích dữ liệu, tái lập mô hình gốc, đến đề xuất mô hình \model{}. \model{} chọn một người nghe làm đại diện cho người sẽ phản hồi mà không dùng thông tin tương lai, biểu diễn mỗi người trong mỗi lượt bằng một token bằng chứng và dùng một bộ mã hóa tập hợp nhỏ (0,71 triệu tham số) để dự đoán. Trên tập kiểm tra được khóa và chỉ đánh giá một lần theo kế hoạch đăng ký trước, \model{} tăng UAR từ 21,40 lên 25,24 và WAR từ 32,52 lên 36,43 so với mô hình gốc được huấn luyện lại; chỉ số chính đã đăng ký trước (UAR sau hiệu chỉnh logit) tăng tương tự nhưng không có ý nghĩa thống kê. Các phân tích có kiểm soát trên dữ liệu phát triển cho thấy khuôn mặt của người nghe và ngữ cảnh của các lượt trước là hai nguồn thông tin quan trọng nhất, mỗi nguồn đóng góp khoảng hai điểm UAR; người nghe được chọn tự động cho kết quả tương đương người phản hồi thật hay một người không nói được chọn ngẫu nhiên, và nhãn vai trò không tạo khác biệt đo được; và lợi thế của mô hình tập trung ở các trường hợp người phản hồi lặp lại cảm xúc của người nói hoặc giữ cảm xúc của chính mình. Báo cáo cũng trình bày các hướng đã thử nhưng không thành công và rút ra rằng giới hạn hiện tại nằm ở việc nhận diện trạng thái cảm xúc của những người tham gia.

\vspace{0.5cm}
\noindent\textbf{\textit{Từ khóa:}} \textit{dự đoán cảm xúc, hội thoại đa phương thức, Hi-EF, người nghe.}
}
